\section{Análisis y Diseño}

El \textbf{Chat Multicliente} puede dividirse fundamentalmente en dos programas independientes:
\textbf{cliente} y \textbf{servidor}. La comunicación entre ambas partes se realiza con \textbf{JSON}
sobre \textbf{TCP} por medio del protocolo establecido en el curso. El protocolo cuenta con \textbf{23}
mensajes establecidos, más un mensaje genérico de respuesta por parte del servidor.

\subsection*{Cliente}

El cliente se diseñó para ser desarrollado en el lenguaje de programación \textbf{Golang}, utilizando 
la librería \textbf{Bubble Tea} para la interfaz de terminal. Esta decisión se tomó de manera previa
al análisis de la especificación del proyecto, pues estaba realmente entusiasmado en utilizar la 
librería para algún proyecto. En consecuencia, \textbf{Gabby} adopta la \textbf{Arquitectura ELM}(\textit{Model-Update-View})
como arquitectura principal. 

La \textbf{Arquitectura ELM} surge en la comunidad del lenguaje de programación funcional \textbf{ELM}
como una manera de estructurar interfaces gráficas reactivas y funcionales. Se compone de tres elementos:

\begin{description}[leftmargin=2.2cm, style=nextline, font=\bfseries]
    \item[Modelo] Representa el estado actual de la aplicación.
    \item[Update] Ciclo que actualiza el Modelo a partir de mensajes recibidos.
    \item[Vista] Una fotografía (\textit{snapshot}) del estado en un momento dado.
\end{description}

Generalmente, las aplicaciones realizadas a partir de la \textbf{Arquitectura ELM} cuentan con un modelo 
raíz que encapsula varios submodelos. Cada submodelo se encarga de la lógica de una parte específica y
delimitada de la interfaz manteniendo su propio estado, su propia función de actualización y su propia vista. 
El modelo raíz orquesta estos submodelos: delega los mensajes correspondientes a cada uno, integra sus vistas parciales 
en la vista general de la aplicación, y coordina la comunicación entre ellos cuando es necesario.

% TODO: Explicar el modelado de los paneles.

\begin{figure}[h]
    \centering
    \includegraphics[width=0.8\textwidth]{landing.pdf}
    \caption{Mockup de la interfaz de terminal del cliente.}
    \label{fig:mockup-cliente}
\end{figure}

\subsection*{Servidor}